\section*{Leitidee}
Die Fourieranalyse beruht auf der Beobachtung, dass sich viele Signale als �berlagerung harmonischer Schwingungen darstellen lassen.  Sinus- und Kosinusfunktionen bilden dabei eine geeignete Basis, auf der komplexe Signale zerlegt werden k�nnen.

Diese Zerlegung l�sst sich als Projektion einer Funktion auf ein Orthonormalsystem harmonischer Schwingungen verstehen.  Dadurch wird ein Signal nicht mehr nur als Funktion im Zeitbereich betrachtet, sondern als Kombination einfacher Schwingungen unterschiedlicher Frequenzen.

In der komplexen Exponentialdarstellung lassen sich diese Basisfunktionen besonders elegant behandeln, da Differentiation, Integration und Superposition zu einfachen algebraischen Operationen werden. Die Eulersche Identit�t
\[
	e^{i\, x}=\cos(x)+i\, \sin(x)
\]
verkn�pft dabei die trigonometrische mit der komplexen Darstellung und macht die zugrunde liegende Struktur sichtbar.
